% sections/05_temporal_analysis.tex
% §5: RQ2 Temporal IC Analysis
% Restructured around IC temporal extension

\section{RQ2: Temporal Inconsistency Evolution}
\label{sec:temporal}

\subsection{Experimental Setup}

We construct a temporal dataset of 24,243 version transition pairs spanning 7,144 \clawhub{} skills (3.39 documented transitions per skill on average).

\noindent\textbf{Exploratory note.} Temporal IC labeling (T4/T5) relies on the same
LSCPS static extraction used for the main IC analysis; temporal labels are therefore
\emph{proxy indicators} derived from version-history metadata, not independently verified.
TDM recall is not included in the main composite evaluation (Table~\ref{tab:full_eval})
because its proxy labels are not independent of the detection features.
The findings here are reported as \emph{exploratory observations} to motivate
future deployment with live version-controlled repositories.
For each unique version pair $(v_i, v_{i+1})$, we compute:
\begin{itemize}[topsep=2pt,itemsep=1pt]
    \item \textbf{TCE} (Eq.~\ref{eq:tce}): Risk-Weighted Mismatch delta across versions;
    \item \textbf{SDS}: fraction of newly introduced capabilities mentioned in the changelog;
    \item \textbf{Family-level drift}: per-family CES deltas (Shell, Credential, Network, File, Agentic).
\end{itemize}
T4/T5 labels are \emph{heuristic risk indicators} from archived text, not verified malicious labels.

\paragraph{Scope note: \clawhub{} only.}
Temporal analysis is \clawhub{}-exclusive; \skillssh{} provides only a static snapshot
without structured version history.
Cross-platform static results are reported in \S\ref{sec:static} (RQ4).

\subsection{Temporal IC Distribution}

\begin{figure}[t]
\centering
\includegraphics[width=\columnwidth]{figures/fig_ic_temporal.pdf}
\caption{Version transition temporal IC labels on \clawhub{} ($M{=}24{,}243$ pairs).
  T6 (confirmed malicious) transitions are shown as a separate category.}
\label{fig:transitions}
\end{figure}

\begin{table}[h]
\centering
\caption{Temporal IC classification of 24,243 version transitions.
  ``Silent expansion'' = TCE$>$0 and SDS$<$0.35.}
\label{tab:temporal_ic}
\scriptsize
\begin{tabular}{@{}lrrp{3.2cm}@{}}
\toprule
\textbf{Label} & \textbf{Count} & \textbf{\%} & \textbf{IC meaning} \\
\midrule
T-Benign         & 19,131 & 78.9\% & No IC degradation \\
T-Silent-Expand  &  1,940 &  8.0\% & TCE$>$0, SDS$<$0.35 (T4) \\
T-Silent-Cred    &  1,136 &  4.7\% & CRED-READ silently expanded \\
T-DocLag         &    623 &  2.6\% & Capability added, description static \\
T-Silent-Minor   &    534 &  2.2\% & Low-risk generic drift \\
T-Silent-Major   &    512 &  2.1\% & High-risk cap injected without disclosure \\
T-CovertAgent    &    327 &  1.3\% & AUTONOMOUS silently added \\
\textbf{T6-Malicious} & \textbf{40} & \textbf{0.2\%} & \textbf{Confirmed malicious injection} \\
\bottomrule
\end{tabular}
\end{table}

\paragraph{Answer to RQ2.}
\textbf{2,575 of 24,243 transitions (10.6\%)} are silent expansions
(TCE$>$0, SDS$<$0.35).
This figure aggregates across all temporal sub-categories in Table~\ref{tab:temporal_ic}
(T-Silent-Expand, T-Silent-Cred, T-DocLag, T-Silent-Minor, T-Silent-Major, and T-CovertAgent)
that satisfy the TCE$>$0, SDS$<$0.35 criterion.
Ratio of silent to disclosed expansions: $2{,}575 : 534 = \textbf{4.8{:}1}$.
The most silently-expanded families:
CRED-READ (39.0\%), SHELL-EXEC (31.6\%), FILE-OPS (17.8\%).
Median changelog for silent expansions: 3 words (``minor bug fix'').

\subsection{Cumulative IC Degradation (T5)}

T5 requires IC-1/IC-2 at the current version
\emph{and} at least one silent expansion in the skill's version history.
We find \textbf{1,376 skills} (7.8\% of all \NskillsDMR{} skills;
19.3\% of the 7,144 skills with version history) satisfy this criterion.
T5 skills are \textbf{4.5$\times$ more likely} to be associated with
a T6 confirmed malicious injection than non-T5 skills
(Fisher's exact test, $p{<}0.01$; OR\,=\,4.5, 95\%\,CI\,[1.8,\,11.2]).
This association suggests T5 may be a useful risk indicator,
but the proxy nature of T4/T5 labels means this cannot be treated
as independent validation of the detection pipeline.

\paragraph{Changelog word count distribution.}
Among silently expanding transitions, the median changelog is 3 words
(e.g., ``minor bug fix,'' ``performance update,'' ``1.0.1'').
Only 4 of 2,575 candidate pairs feature changelogs exceeding 50 words.

\paragraph{Temporal pattern.}
Skills do not expand capabilities linearly.
A ``plateau-and-jump'' pattern emerges:
78\% of aggregate CES expansions occur within a single version transition
rather than gradually over time.
This concentration is consistent with abrupt capability change,
though not sufficient on its own to establish malicious intent.

\subsection{Version Lifecycle IC Evolution}
\label{sec:lifecycle}

A natural question is whether \emph{developers self-correct} over time:
do skills become more transparent in later versions?
We group transitions by the target version number $v$ into five cohorts
and compute per-cohort IC-0 rate (consistent disclosure) and mean CR.

\begin{table}[h]
\centering
\small
\setlength{\tabcolsep}{4pt}
\caption{IC disclosure quality across version lifecycle cohorts (24,243 transitions).
  All cohorts are sub-sampled by version number; v1 ($N{=}21$) and v2--3 ($N{=}54$) have
  small $N$ and should be interpreted cautiously.}
\label{tab:lifecycle_ic}
\begin{tabular}{@{}lrrrrr@{}}
\toprule
\textbf{Cohort} & \textbf{N} & \textbf{IC-0\%} & \textbf{IC-1\%} & \textbf{CR} & \textbf{Silent\%} \\
\midrule
v1   (Initial)     & 21  & 4.8  & 61.9 & 0.221 & 0.0 \\
v2--3 (Early)      & 54  & 7.4  & 51.9 & 0.252 & 5.6 \\
v4--7 (Mid)        & 100 & 10.0 & 45.0 & 0.274 & \textbf{11.0} \\
v8--15 (Mature)    & 191 & 11.0 & 55.5 & 0.313 & 0.5 \\
v16+  (Long-lived) & 134 & 20.1 & 56.0 & 0.357 & 0.0 \\
\bottomrule
\end{tabular}
\end{table}

\begin{figure}[t]
\centering
\includegraphics[width=\columnwidth]{figures/fig_temporal_ic_evolution.pdf}
\caption{IC distribution and mean CR by version cohort, with silent expansion rate.
Silent expansions peak in the mid-lifecycle (v4--7), the period of active development
but limited security scrutiny.}
\label{fig:lifecycle}
\end{figure}

\paragraph{Key findings.}
Mean CR improves slightly ($0.221 \to 0.357$) as skills age,
suggesting limited but real developer self-correction over time.
However, \textbf{IC-1 (Full Omission) remains the dominant IC type in all cohorts},
including long-lived skills ($>$v16): 56\% of transitions in that cohort
still exhibit complete under-disclosure.
This pattern is \textbf{consistent with} IC violations being a structural ecosystem property
rather than a developer-learning-curve artifact---though the cohort sizes
(especially v1--v3) are small and caution is warranted in generalizing.

Notably, silent capability expansions \textbf{peak at v4--7 (11\%)},
the mid-lifecycle window where skills are actively developed
but not yet subject to mature security scrutiny or active user review.
This is the \emph{highest-risk window} for undetected capability injection,
providing a deployment heuristic: skills in this version band should receive
higher-priority review from TrustDrift's screening queue.
